% Generated draft: requires mathematical and source review.
\documentclass[11pt]{article}
\usepackage[margin=25mm]{geometry}
\usepackage{fontspec}
\setmainfont{Latin Modern Roman}
\usepackage{amsmath,amssymb,mathtools,unicode-math}
\setmathfont{Latin Modern Math}
\usepackage{enumitem,needspace,xurl,hyperref}
\setlength{\parindent}{0pt}
\setlength{\parskip}{5pt}
\setlength{\emergencystretch}{3em}
\newcommand{\N}{\mathbb{N}}
\newcommand{\Z}{\mathbb{Z}}
\newcommand{\Q}{\mathbb{Q}}
\newcommand{\R}{\mathbb{R}}
\newcommand{\C}{\mathbb{C}}
\newcommand{\F}{\mathbb{F}}
\newcommand{\A}{\mathbb{A}}
\newcommand{\PP}{\mathbb{P}}
\newcommand{\E}{\mathbb{E}}
\newcommand{\eps}{\varepsilon}
\newcommand{\dd}{\,\mathrm{d}}
\begin{document}
\section*{UnsolvedMath ID: 9000001}
\textbf{Source problem code:} catalogue entry
\subsection{Definitions and mathematical statement}
Let \(\Omega\) be the source's compact convex domain, let \(F_\Omega\) be its assigned function, and let \(C_\Omega\) be the corner locus. For a piecewise-affine function, a corner locus is the set where it is not locally affine, with vertices supplied by the associated polyhedral structure. The source-specific function and that structure are not defined in this extraction. The stated coordinate data are
\[
c(\Omega)=\bigl(F_\Omega(v)\bigr)_{v\in V(C_\Omega)}.
\]
No finiteness, labeling, or fixed combinatorial type of this vertex set is assumed.

The source says these are complete coordinates and asks for an explicit relation between
\[
c(\Omega)\quad\text{and}\quad c(\Omega^*).
\]
Equivalently, in a specified coordinate/moduli framework, describe the induced duality operation \(\mathcal D\) with
\[
c(\Omega^*)=\mathcal D(c(\Omega)).
\]
An abstract definition by reconstructing the domain and then dualizing would not by itself provide the explicit coordinate relation being sought. No coordinatewise pairing of vertices is presumed; duality may change the corner-locus combinatorics.

The dual-domain convention is missing. For example, if the intended planar convention is polar duality, one needs a chosen origin with \(0\in\operatorname{int}\Omega\) and defines
\[
\Omega^*=\{y\in\R^2:\langle x,y\rangle\le1\text{ for all }x\in\Omega\}.
\]
This requires the origin/normalization and is different from other projective or Legendre-dual conventions. The meaning of complete coordinates, including translation or other equivalences and vertex labeling, must be compatible with whichever duality is intended. Neither that polar convention nor a formula for \(F_\Omega\) is silently substituted for the absent source definitions.

The supplied research summary reports later tropical-caustic moduli progress toward the coordinate programme but says the duality relation remains open and not fully formalized. No explicit transformation or complete follow-up citation is supplied.
\subsection{Short English statement}
Express the dual convex domain's corner-locus vertex values in terms of the original domain's complete coordinate data. The function, indexing/normalization, and duality convention must first be fixed.
\subsection{Sources}
\begin{itemize}
\item Local UnsolvedMath registry, permanent ID \(9000001\), exact source code \texttt{catalogue entry}, ``Coordinates on convex domains.'' Proposers Nikita Kalinin and Mikhail Shkolnikov.
\item Kalinin and Shkolnikov, \emph{The Number pi and a Summation by SL(2,Z)} (2018), section 4, Questions direction 1; \url{https://amj.math.stonybrook.edu/html-articles/Files-2015-2024/17-75/index.html}. HTML-text extraction; recorded access 2026-07-29.
\item Supplied PARTIAL-PROGRESS summary mentions a follow-up tropical-caustic moduli paper without a title, publication details, or a coordinate-duality theorem.
\end{itemize}
\end{document}
